\clearpage
\oldpage[IV-1]{251}
\phantomsection
\addcontentsline{toc}{section}{Bibliografía (continuación)}
\section*{BIBLIOGRAFÍA \textit{(continuación)}}

\begingroup
\small
\setlength{\parindent}{0pt}
\setlength{\parskip}{3pt}
[30] M.~\textsc{Nagata}, Local rings, \emph{Interscience Tracts}, 13 (1962),
Interscience, New York.

[31] M.~\textsc{Nagata}, A Jacobian criterion of simple points,
\emph{Ill. Journ. of Math.}, 1 (1957), p.~427-432.

[32] A.~\textsc{Néron}, Modèles minimaux des variétés abéliennes sur les corps
locaux et globaux, \emph{Publ. Inst. Hautes Études Sci.}, n\textsuperscript{o}
21 (1964).

\clearpage
\oldpage[IV-1]{252}
\phantomsection
\addcontentsline{toc}{section}{Índice de notaciones}
\section*{ÍNDICE DE NOTACIONES}

$\operatorname{dimc}(X)$, $\dim(X)$ ($X$ espacio topológico): \textbf{0}, 14.1.2.\par
$\dim_x(X)$ ($X$ espacio topológico, $x$ punto de $X$): \textbf{0}, 14.1.2.\par
$\operatorname{codim}(Y,X)$ ($X$ espacio topológico, $Y$ subconjunto cerrado de $X$): \textbf{0}, 14.2.1.\par
$\operatorname{codim}_x(Y,X)$ ($X$ espacio topológico, $Y$ subconjunto cerrado de $X$, $x$ punto de $X$): \textbf{0}, 14.2.4.\par
$N[T_1,\ldots,T_n]$ ($N$ $A$-módulo): \textbf{0}, 15.1.1.\par
$\dim(A)$ ($A$ anillo): \textbf{0}, 16.1.1.\par
$\operatorname{ht}(\mathfrak a)$ ($\mathfrak a$ ideal): \textbf{0}, 16.1.3.\par
$\dim_A(M)$, $\dim(M)$ ($M$ $A$-módulo): \textbf{0}, 16.1.7.\par
$P_{\mathfrak q}(M,n)$ ($M$ módulo de tipo finito sobre un anillo semilocal noetheriano $A$, $\mathfrak q$ ideal de definición de $A$): \textbf{0}, 16.2.1.\par
$\operatorname{prof}_A(M)$, $\operatorname{prof}(M)$ ($M$ $A$-módulo): \textbf{0}, 16.4.5.\par
$\operatorname{coprof}_A(M)$, $\operatorname{coprof}(M)$ ($M$ $A$-módulo): \textbf{0}, 16.4.9.\par
$\operatorname{dim.proj}_A(M)$, $\operatorname{dim.proj}(M)$,
$\operatorname{dim.inj}_A(M)$, $\operatorname{dim.inj}(M)$ ($M$ $A$-módulo): \textbf{0}, 17.2.1.\par
$\operatorname{dim.coh}(A)$ ($A$ anillo): \textbf{0}, 17.2.8.\par
$\operatorname{dim.proj}(\mathscr F)$, $\operatorname{dim.inj}(\mathscr F)$
($\mathscr F$ Módulo sobre un espacio anillado): \textbf{0}, 17.2.14.\par
$\operatorname{dim.coh}(X)$ ($X$ espacio anillado): \textbf{0}, 17.2.14.\par
$E\times_B F$ ($B$, $E$, $F$ $A$-anillos): \textbf{0}, 18.1.2.\par
$D_B(L)$ ($B$ anillo, $L$ $B$-bimódulo): \textbf{0}, 18.2.3.\par
$Y\oplus_X Z$ ($X$, $Y$, $Z$ $A$-bimódulos): \textbf{0}, 18.2.7.\par
$\operatorname{Exan}_A(B,L)$ ($B$ $A$-anillo, $L$ $B$-bimódulo): \textbf{0}, 18.3.4.\par
$\operatorname{Exan}_{A/A'}(B,L)$: \textbf{0}, 18.3.7.\par
$\operatorname{Exal}_A(B,L)$, $\operatorname{Exal}_{A/A'}(B,L)$ ($A$ anillo conmutativo, $B$ $A$-álgebra, $L$ $B$-bimódulo): \textbf{0}, 18.4.1.\par
$\operatorname{Exalcom}_A(B,L)$, $\operatorname{Exalcom}_{A/A'}(B,L)$ ($A$ anillo conmutativo, $B$ $A$-álgebra conmutativa, $L$ $B$-módulo): \textbf{0}, 18.4.2.\par
$\mathrm H_A^i(B,L)$, $\mathrm H_A^2(B,L)^s$: \textbf{0}, 18.4.3.\par
$\operatorname{Exantop}_A(B,L)$, $\operatorname{Exaltop}_A(B,L)$,
$\operatorname{Exalcotop}_A(B,L)$: \textbf{0}, 18.5.1.\par
$\operatorname{Exantop}_{A/A'}(B,L)$, $\operatorname{Exaltop}_{A/A'}(B,L)$,
$\operatorname{Exalcotop}_{A/A'}(B,L)$: \textbf{0}, 18.5.2.\par
$\operatorname{Hom.\ cont}_C(M,N)$ ($M$, $N$ $C$-bimódulos topológicos): \textbf{0}, 18.5.3.\par
$\operatorname{D\acute{e}r}_A(B,L)$ ($B$ $A$-anillo, $L$ $B$-bimódulo): \textbf{0}, 20.1.2.\par
$\operatorname{D\acute{e}r.\ cont}_A(B,L)$: \textbf{0}, 20.3.1.\par
$p_{B/A}$, $\mathfrak J_{B/A}$ ($B$ $A$-álgebra topológica conmutativa): \textbf{0}, 20.4.1.\par
$P^1_{B/A}$, $\varepsilon_{B/A}$ ($B$ $A$-álgebra topológica conmutativa): \textbf{0}, 20.4.2.\par
$\Omega^1_{B/A}$, $\widehat{\Omega}^1_{B/A}$ ($B$ $A$-álgebra topológica conmutativa): \textbf{0}, 20.4.3.\par
$\Omega^1_B$ ($B$ anillo topológico): \textbf{0}, 20.4.3.\par
$d_{B/A}$, $d(x)$, $dx$, $d_B$: \textbf{0}, 20.4.6.\par
$\Upsilon_{C/B/A}$, $\Upsilon_{C/B}$ ($A$, $B$, $C$ anillos): \textbf{0}, 20.6.1.\par
$\chi_E$ ($E$ $B$-extensión $A$-trivial de $C$ por $L$): \textbf{0}, 20.6.8.\par

\clearpage
\oldpage[IV-1]{253}
$\Upsilon^C_{B/A/\Lambda}$, $\Upsilon^C_{B/A}$ ($\Lambda$, $A$, $B$, $C$ anillos): \textbf{0}, 20.6.14.\par
$\chi_{B/A}$, $\chi_B$ ($B$ $A$-álgebra): \textbf{0}, 20.6.24.\par
$\chi_E$ ($E$ $B$-álgebra topológica): \textbf{0}, 20.7.10,\par
$F_A$, $A^p$, $A^{(p)}$, $E^{(p)}$ ($A$ anillo de característica $p>0$, $E$ $A$-módulo): \textbf{0}, 21.1.4.\par
$\pi_{B/A}$, $\Theta_{B/A}$, $\Xi_{B/A}$ ($A$, $B$ anillos de característica $p>0$): \textbf{0}, 21.3.2.\par
$\Delta(L/K,k/k_0)$, $d(L/K,k/k_0)$, $\Delta(L/K,k)$, $d(L/K,k)$ ($k_0$, $k$, $K$, $L$ cuerpos): \textbf{0}, 21.6.1.\par
$\chi'_{B/A}$ ($A$, $B$ anillos de característica $p>0$): \textbf{0}, 22.4.6.\par
$\mathscr T^{\mathrm{cons}}$, $X^{\mathrm{cons}}$ ($X$ preesquema, $\mathscr T$ topología de $X$): \textbf{IV}, 1.9.12.\par
$f^{\mathrm{cons}}$ ($f$ morfismo de preesquemas): \textbf{IV}, 1.9.13.\par

\clearpage
\oldpage[IV-1]{254}
\phantomsection
\addcontentsline{toc}{section}{Índice terminológico}
\section*{ÍNDICE TERMINOLÓGICO}

Álgebra aumentada de las partes principales de orden 1: \textbf{0}, 20.4.2.\par
$\mathbf{A}$-álgebra de Cohen: \textbf{0}, 19.8.1.\par
Álgebra de presentación finita: \textbf{IV}, 1.4.1.\par
$A$-álgebra esencialmente de tipo finito: \textbf{IV}, 1.3.8.\par
$A$-anillo: \textbf{0}, 18.1.1.\par
$A$-anillo aumentado sobre $B$: \textbf{0}, 18.1.4.\par
Anillo biequidimensional: \textbf{0}, 16.1.4.\par
Anillo catenario: \textbf{0}, 16.1.4.\par
Anillo de característica $p$: \textbf{0}, 21.1.1.\par
Anillo de Cohen: \textbf{0}, 19.8.4.\par
Anillo de Cohen-Macaulay: \textbf{0}, 16.5.1 y 16.5.13.\par
Anillo equicodimensional, equidimensional: \textbf{0}, 16.1.4.\par
Anillo local geométricamente unirrama: \textbf{0}, 23.1.7.\par
Anillo japonés: \textbf{0}, 23.1.1.\par
Anillo local primo, anillo local completo primo: \textbf{0}, 19.8.3.\par
Anillo regular: \textbf{0}, 17.1.1 y 17.3.6.\par
Anillo local unirrama: \textbf{0}, 23.1.7.\par
Anillo universalmente japonés: \textbf{0}, 23.1.1.\par
Aplicación abierta en un punto: \textbf{IV}, 1.10.2.\par
Aumentación: \textbf{0}, 18.1.4.\par

Bimorfismo formal: \textbf{0}, 19.1.2.\par

Cadena en un conjunto ordenado, cadena de subconjuntos cerrados irreducibles de un espacio topológico: \textbf{0}, 14.1.1.\par
Cadena saturada: \textbf{0}, 14.3.1.\par
Codimensión combinatoria, codimensión de un subconjunto cerrado de un espacio topológico: \textbf{0}, 14.2.1.\par
Codimensión en un punto de un subconjunto cerrado de un espacio topológico: \textbf{0}, 14.2.4.\par
Condición de las cadenas: \textbf{0}, 14.3.2.\par
Coprofundidad de un módulo sobre un anillo local noetheriano: \textbf{0}, 16.4.9.\par
Cuerpo $k_0$-admisible, cuerpo admisible para una extensión: \textbf{0}, 21.6.1.\par
Cuerpo de multiplicidad radicial finita sobre un subcuerpo: \textbf{0}, 19.6.6.\par

Defecto de $k_0$-admisibilidad de una extensión: \textbf{0}, 21.6.1.\par
Derivación, $A$-derivación: \textbf{0}, 20.1.2.\par
Diferenciales, 1-diferenciales, diferenciales absolutas: \textbf{0}, 20.4.3.\par
Diferencial exterior: \textbf{0}, 20.4.6.\par
Dimensión combinatoria, dimensión de un espacio topológico: \textbf{0}, 14.1.2.\par
Dimensión de un espacio topológico en un punto: \textbf{0}, 14.1.2.\par
Dimensión de Krull, dimensión de un anillo: \textbf{0}, 16.1.1.\par
Dimensión de un módulo: \textbf{0}, 16.1.7.\par
Dimensión cohomológica, dimensión cohomológica global de un anillo: \textbf{0}, 17.2.8.\par
Dimensión cohomológica de un espacio anillado en un punto, dimensión cohomológica de un espacio anillado: \textbf{0}, 17.2.14.\par

\clearpage
\oldpage[IV-1]{255}
Dimensión inyectiva, dimensión proyectiva de un módulo: \textbf{0}, 17.2.1.\par
Dimensión inyectiva en un punto, dimensión inyectiva de un haz: \textbf{0}, 17.2.14.\par
Dimensión proyectiva en un punto, dimensión proyectiva de un haz: \textbf{0}, 17.2.14.\par

Elemento de $A$ no divisor de 0 en un $A$-módulo: \textbf{0}, 15.1.1.\par
Elemento $M$-regular, elemento $M$-cuasirregular: \textbf{0}, 15.1.4.\par
Conjunto ind-construible, conjunto pro-construible: \textbf{IV}, 1.9.4.\par
Epimorfismo formal: \textbf{0}, 19.1.2.\par
$A$-equivalencia de $A$-extensiones: \textbf{0}, 18.2.2.\par
Espacio biequidimensional: \textbf{0}, 14.3.3.\par
Espacio catenario: \textbf{0}, 14.3.2.\par
Espacio equicodimensional: \textbf{0}, 14.2.1.\par
Espacio equidimensional: \textbf{0}, 14.1.3.\par
$A$-extensión de un $A$-anillo $B$ por un $B$-bimódulo: \textbf{0}, 18.2.2.\par
Extensión deducida de otra por un homomorfismo de bimódulos: \textbf{0}, 18.2.8.\par
Extensión de Hochschild: \textbf{0}, 18.4.3.\par
$B$-extensión $A$-trivial: \textbf{0}, 18.3.7.\par
Extensión trivial tipo: \textbf{0}, 18.2.3.\par
$A$-extensiones $A$-equivalentes: \textbf{0}, 18.2.2.\par

Familia $p$-libre, familia absolutamente $p$-libre: \textbf{0}, 21.1.9.\par
Formalmente étale (álgebra): \textbf{0}, 19.10.2.\par
Formalmente invertible a la derecha (homomorfismo): \textbf{0}, 19.1.15.\par
Formalmente invertible a la izquierda (homomorfismo): \textbf{0}, 19.1.5.\par
Formalmente no ramificada (álgebra): \textbf{0}, 19.10.2.\par
Formalmente proyectivo (módulo): \textbf{0}, 19.2.1.\par
Formalmente suave (álgebra): \textbf{0}, 19.3.1.\par
Formalmente suave con relación a un anillo (álgebra): \textbf{0}, 19.9.1.\par

Geométricamente regular sobre un cuerpo (anillo local): \textbf{0}, 19.6.5.\par
Grupo de las clases de $A$-extensiones: \textbf{0}, 18.3.4.\par

Altura de un ideal: \textbf{0}, 16.1.3.\par
$A$-homomorfismo de $A$-anillos: \textbf{0}, 18.1.1.\par
Homomorfismo característico de una $B$-extensión $A$-trivial: \textbf{0}, 20.6.8.\par
Homomorfismo característico de una $A$-álgebra con relación a un anillo y a un ideal: \textbf{0}, 20.6.24.\par
Homomorfismo característico de un álgebra de característica $p>0$: \textbf{0}, 22.4.6.\par

Ideal de aumentación: \textbf{0}, 18.1.4.\par
Imagen inversa de un anillo aumentado: \textbf{0}, 18.1.5.\par
Imagen inversa de una $A$-extensión: \textbf{0}, 18.2.5.\par

Longitud de una cadena: \textbf{0}, 14.1.1.\par

Módulo biequidimensional, catenario, equicodimensional, equidimensional: \textbf{0}, 16.1.7.\par
Módulo de Cohen-Macaulay: \textbf{0}, 16.5.1 y 16.5.13.\par
Módulo de imperfección: \textbf{0}, 20.6.1.\par
Monomorfismo formal: \textbf{0}, 19.1.2.\par
Morfismo de $A$-extensiones: \textbf{0}, 18.2.4.\par
Morfismo de presentación finita en un punto, morfismo localmente de presentación finita: \textbf{IV}, 1.4.2.\par
Morfismo de presentación finita: \textbf{IV}, 1.6.1.\par
Morfismo cuasiseparado: \textbf{IV}, 1.2.1.\par

$p$-anillo de Cohen: \textbf{0}, 19.8.4.\par
$p$-base, $p$-base absoluta: \textbf{0}, 21.1.9.\par
Punto maximal de un preesquema: \textbf{IV}, 1.1.4.\par

\clearpage
\oldpage[IV-1]{256}
Polinomio de Hilbert-Samuel: \textbf{0}, 16.2.1.\par
Preesquema de presentación finita sobre otro: \textbf{IV}, 1.6.1.\par
Preesquema cuasiseparado: \textbf{IV}, 1.2.1.\par
Producto fibrado de $A$-anillos: \textbf{0}, 18.1.2.\par
Profundidad de un módulo: \textbf{0}, 16.4.5.\par

Suma amalgamada de bimódulos: \textbf{0}, 18.2.7.\par
Estrictamente formalmente proyectivo (módulo): \textbf{0}, 19.2.3.\par
Sucesión $M$-regular, sucesión $M$-cuasirregular: \textbf{0}, 15.1.7.\par
Sucesión $\mathscr F$-regular, sucesión $\mathscr F$-cuasirregular: \textbf{0}, 15.2.2.\par
Sistema de parámetros para un módulo sobre un anillo semilocal noetheriano: \textbf{0}, 16.3.6.\par
Sistema de $p$-generadores, sistema absoluto de $p$-generadores: \textbf{0}, 21.1.9.\par
Sistema fundamental de ideales abiertos, de submódulos abiertos: \textbf{0}, 19.0.2.\par
Sistema regular de parámetros de un anillo local regular: \textbf{0}, 17.1.6.\par

Topología construible sobre un preesquema: \textbf{IV}, 1.9.12.\par
Topología sobre un $A$-módulo deducida de la topología de $A$: \textbf{0}, 19.0.2.\par
Trivial ($A$-anillo aumentado): \textbf{0}, 18.1.4.\par
$A$-trivial (extensión): \textbf{0}, 18.2.3 y 18.3.7.\par

\clearpage
\oldpage[IV-1]{257}
\phantomsection
\addcontentsline{toc}{section}{Índice general}
\section*{ÍNDICE GENERAL}

\hfill\textsc{Páginas}

\newcommand{\ivonecontentsline}[3]{\noindent\makebox[17mm][l]{#1}\parbox[t]{0.77\linewidth}{#2\dotfill}\makebox[12mm][r]{#3}\par}
\ivonecontentsline{}{\textsc{Capítulo 0.} --- \textbf{Preliminares} \textit{(continuación)}}{5}
\ivonecontentsline{\S~14.}{Dimensión combinatoria de un espacio topológico}{6}
\ivonecontentsline{14.1.}{Dimensión combinatoria de un espacio topológico}{6}
\ivonecontentsline{14.2.}{Codimensión de un subconjunto cerrado}{8}
\ivonecontentsline{14.3.}{La condición de las cadenas}{10}
\ivonecontentsline{\S~15.}{Sucesiones $M$-regulares y sucesiones $\mathscr F$-regulares}{12}
\ivonecontentsline{15.1.}{Sucesiones $M$-regulares y sucesiones $M$-cuasirregulares}{12}
\ivonecontentsline{15.2.}{Sucesiones $\mathscr F$-regulares}{20}
\ivonecontentsline{\S~16.}{Dimensión y profundidad en los anillos locales noetherianos}{22}
\ivonecontentsline{16.1.}{Dimensión de un anillo}{22}
\ivonecontentsline{16.2.}{Dimensión de un anillo semilocal noetheriano}{25}
\ivonecontentsline{16.3.}{Sistemas de parámetros en un anillo local noetheriano}{28}
\ivonecontentsline{16.4.}{Profundidad y coprofundidad}{32}
\ivonecontentsline{16.5.}{Módulos de Cohen-Macaulay}{36}
\ivonecontentsline{\S~17.}{Anillos regulares}{39}
\ivonecontentsline{17.1.}{Definición de los anillos regulares}{39}
\ivonecontentsline{17.2.}{Recordatorios sobre la dimensión proyectiva y la dimensión inyectiva de los módulos}{42}
\ivonecontentsline{17.3.}{Teoría cohomológica de los anillos regulares}{46}
\ivonecontentsline{\S~18.}{Complementos sobre las extensiones de álgebras}{51}
\ivonecontentsline{18.1.}{Imágenes inversas de anillos aumentados}{51}
\ivonecontentsline{18.2.}{Extensiones de un anillo por un bimódulo}{54}
\ivonecontentsline{18.3.}{El grupo de las clases de $A$-extensiones}{59}
\ivonecontentsline{18.4.}{Extensiones de álgebras}{64}
\ivonecontentsline{18.5.}{Caso de los anillos topológicos}{66}
\ivonecontentsline{\S~19.}{Álgebras formalmente suaves y anillos de Cohen}{69}
\ivonecontentsline{19.0.}{Introducción}{69}
\ivonecontentsline{19.1.}{Epimorfismos y monomorfismos formales}{71}
\ivonecontentsline{19.2.}{Módulos formalmente proyectivos}{78}

\clearpage
\oldpage[IV-1]{258}
\hfill\textsc{Páginas}

\ivonecontentsline{19.3.}{Álgebras formalmente suaves}{79}
\ivonecontentsline{19.4.}{Primeros criterios de suavidad formal}{86}
\ivonecontentsline{19.5.}{Suavidad formal y anillos graduados asociados}{90}
\ivonecontentsline{19.6.}{Caso de las álgebras sobre un cuerpo}{100}
\ivonecontentsline{19.7.}{Caso de los homomorfismos locales: teoremas de existencia y de unicidad}{104}
\ivonecontentsline{19.8.}{Álgebras de Cohen y $p$-anillos de Cohen; aplicación a la estructura de los anillos locales completos}{109}
\ivonecontentsline{19.9.}{Álgebras relativamente formalmente suaves}{114}
\ivonecontentsline{19.10.}{Álgebras formalmente no ramificadas y álgebras formalmente étales}{115}
\ivonecontentsline{\S~20.}{Derivaciones y diferenciales}{116}
\ivonecontentsline{20.1.}{Derivaciones y extensiones de álgebras}{117}
\ivonecontentsline{20.2.}{Propiedades funtoriales de las derivaciones}{119}
\ivonecontentsline{20.3.}{Derivaciones continuas en los anillos topológicos}{121}
\ivonecontentsline{20.4.}{Partes principales y diferenciales}{123}
\ivonecontentsline{20.5.}{Propiedades funtoriales fundamentales de $\Omega^1_{B/A}$}{128}
\ivonecontentsline{20.6.}{Módulos de imperfección y homomorfismos característicos}{136}
\ivonecontentsline{20.7.}{Generalizaciones a los anillos topológicos}{147}
\ivonecontentsline{\S~21.}{Diferenciales en los anillos de característica $p$}{153}
\ivonecontentsline{21.1.}{Sistemas de $p$-generadores y $p$-bases}{154}
\ivonecontentsline{21.2.}{$p$-bases y suavidad formal}{157}
\ivonecontentsline{21.3.}{$p$-bases y módulos de imperfección}{160}
\ivonecontentsline{21.4.}{Caso de las extensiones de cuerpos}{162}
\ivonecontentsline{21.5.}{Aplicación: criterios de separabilidad}{164}
\ivonecontentsline{21.6.}{Cuerpos admisibles para una extensión}{167}
\ivonecontentsline{21.7.}{La igualdad de Cartier}{169}
\ivonecontentsline{21.8.}{Criterios de admisibilidad}{171}
\ivonecontentsline{21.9.}{Módulos de diferenciales completados en los anillos de series formales}{176}
\ivonecontentsline{\S~22.}{Criterios diferenciales de suavidad formal y de regularidad}{182}
\ivonecontentsline{22.1.}{Levantamiento de la suavidad formal}{183}
\ivonecontentsline{22.2.}{Caracterización diferencial de las álgebras locales formalmente suaves sobre un cuerpo}{186}
\ivonecontentsline{22.3.}{Aplicación a las relaciones entre ciertos anillos locales y sus completados}{191}
\ivonecontentsline{22.4.}{Resultados preliminares sobre las extensiones finitas de anillos locales cuyo ideal maximal es de cuadrado nulo}{193}

\clearpage
\oldpage[IV-1]{259}
\hfill\textsc{Páginas}

\ivonecontentsline{22.5.}{Álgebras geométricamente regulares y álgebras formalmente suaves}{201}
\ivonecontentsline{22.6.}{Criterio jacobiano de Zariski}{205}
\ivonecontentsline{22.7.}{El criterio jacobiano de Nagata}{209}
\ivonecontentsline{\S~23.}{Anillos japoneses}{213}
\ivonecontentsline{23.1.}{Anillos japoneses}{213}
\ivonecontentsline{23.2.}{Clausura integral de un anillo local noetheriano íntegro}{217}
\ivonecontentsline{}{\textsc{Capítulo IV.} --- \textbf{Estudio local de los esquemas y de los morfismos de esquemas}}{222}
\ivonecontentsline{\S~1.}{Condiciones de finitud relativas. Conjuntos construibles en los preesquemas}{224}
\ivonecontentsline{1.1.}{Morfismos cuasicompactos}{224}
\ivonecontentsline{1.2.}{Morfismos cuasiseparados}{226}
\ivonecontentsline{1.3.}{Morfismos localmente de tipo finito}{228}
\ivonecontentsline{1.4.}{Morfismos localmente de presentación finita}{230}
\ivonecontentsline{1.5.}{Morfismos de tipo finito}{233}
\ivonecontentsline{1.6.}{Morfismos de presentación finita}{234}
\ivonecontentsline{1.7.}{Mejora de resultados anteriores}{236}
\ivonecontentsline{1.8.}{Morfismos de presentación finita y conjuntos construibles}{238}
\ivonecontentsline{1.9.}{Conjuntos pro-construibles e ind-construibles}{241}
\ivonecontentsline{1.10.}{Aplicaciones a los morfismos abiertos}{249}
\ivonecontentsline{}{\textsc{Bibliografía} \textit{(continuación)}}{251}
\ivonecontentsline{}{\textsc{Índice de notaciones}}{252}
\ivonecontentsline{}{\textsc{Índice terminológico}}{254}

\vspace{2em}
\begin{flushright}
\emph{Recibido el 15 de marzo de 1963.}
\end{flushright}
\let\ivonecontentsline\relax
\endgroup
